\chapter{Modele recapitulative}

\section{Model 1}

\textbf{Acesta este examenul dat de Prof.\ Olteanu la ACS, sesiunea ianuarie 2018.}

\textbf{Numărul 1}

1. Să se afle valorile extreme ale funcției:
\[
  f(x,y) = x^2 + y^2 - 3x - 2y + 1
\]
pe mulțimea $K: x^2 + y^2 \leq 1$.

\vspace{1cm}

2. Calculați $\ds\int_0^\infty \frac{dx}{x^4 + 1}$.

\vspace{1cm}

3. Să se calculeze volumul corpului mărginit de suprafețele:
\[
  \begin{cases}
    z &= x^2 + y^2 \\
    2 &= x^2 + y^2 + z^2.
  \end{cases}
\]

\vspace{1cm}

4. Să se calculeze fluxul cîmpului $\vec{V} = (x + z^2, y + z^2, -2z)$
prin suprafața $x^2 + y^2 + z^2 = 1, z \geq 0$.

\vspace{1cm}

5. Să se calculeze aria suprafeței:
\[
  \Sigma = \{ (x, y, z) \mid x^2 + y^2 + z^2 = R^2, z \geq 0 \} \cap %
  \{ (x, y, z) \mid x^2 + y^2 \leq Ry \}.
\]
(Aria lui Viviani)

\vspace{1cm}

\textbf{Numărul 2}

1. Să se afle valorile extreme ale funcției:
\[
  f(x, y) = x \cdot y, \text{ pe } K: 2x^2 + y^2 \leq 1.
\]

\vspace{1cm}

2. Calculați $\ds\int_0^\infty \frac{dx}{x^3 + 1}$.

\vspace{1cm}

3. Să se calculeze volumul corpului obținut prin intersecția suprafețelor:
\[
  \begin{cases}
    2 &= x^2 + y^2 + z^2 \\
    z^2 &= x^2 + y^2, z \geq 0
  \end{cases}
\]

\vspace{1cm}

4. Să se calculeze aria suprafeței $2 - z = x^2 + y^2, z \in [0, 1]$.

\vspace{1cm}

5. Calculați fluxul cîmpului vectorial $\vec{V} = \frac{q}{4 \pi} \cdot \frac{\vec{r}}{r^3}$
prin:
\begin{enumerate}[(a)]
\item Sfera centrată în origine și cu rază $R > 0$;
\item Orice suprafață $\Sigma$ închisă, care nu conține originea.
\end{enumerate}
Observație: $\vec{r} = (x, y, z)$ și $r = ||\vec{r}|| = \sqrt{x^2 + y^2 + z^2}$.

(Legea lui Gauss)

%%%%%%%%%%%%%%%%%%%%%%%%%%%%%%%%%%%%%%%%%%%%%%%%%%%%%%%%%%%%%%%%%%%%%%

\section{Model 2}

\indent\indent 1. Calculați, folosind beta și gamma,
$ \ds\int_0^\infty \frac{\sqrt[4]{x}}{(x + 1)^2} dx $.

\vspace{1cm}

2. Calculați integrala:
\[
  \iint_D (1 + \sqrt{x^2 + y^2}) dxdy, \quad D: x^2 + y^2 \leq 2, y > 0.
\]

\vspace{1cm}

3. Calculați volumul corpului mărginit de suprafețele:
\[
  \begin{cases}
    z &= x^2 + y^2 \\
    x^2 + y^2 + z^2 &= 6, z > 0.
  \end{cases}
\]

\vspace{1cm}

4. Calculația aria suprafeței determinate de:
\[
  z = 5 - x^2 - y^2, \quad z \in [1, 4].
\]

\vspace{1cm}

5. Determinați valorile extreme pentru funcția:
\[
  f : D \to \RR, \quad f(x, y) = x^3 + 8y^3 - 2xy,
\]
unde $ D = \{(x, y) \in \RR^2 \mid x, y, \geq 0, y + 2x \leq 2 \} $.

%%%%%%%%%%%%%%%%%%%%%%%%%%%%%%%%%%%%%%%%%%%%%%%%%%%%%%%%%%%%%%%%%%%%%%

\section{Model 3}

\indent\indent 1. Determinați valorile extreme pentru funcția:
\[
  f : \RR^2 \to \RR, \quad f(x, y) = x^4 + y^3 - 4x^3 - 3y^2 + 3y,
\]
unde $ D = \{(x, y) \in \RR \mid x^2 + y^2 \leq 4 \} $.

\vspace{1cm}

2. Calculați $ \ds\int_0^\infty \dfrac{x^2}{(1 + x^4)^2} dx $.

\vspace{1cm}

3. Calculați volumul mulțimii mărginite de suprafețele:
\begin{align*}
  S_1 &: x^2 + y^2 + z^2 = 6z, z \geq 3 \\
  S_2 &: 2(x^2 + y^2) = z^2.
\end{align*}

\vspace{1cm}

4. Fie suprafața $ S = \{ x^2 + y^2 + z^2 = 5, z \leq 0 \} $.
Desenați suprafața și calculați:
\[
  \iint_S (2x + 5x) dy \wedge dz + 2ydz \wedge dx + (2z - 5x) dx \wedge dy.
\]


%%% Local Variables:
%%% mode: latex
%%% TeX-master: "../m1-18-19"
%%% End:
